% topic: algorithms
\question กำหนดขั้นตอนการรวมกลุ่มตัวเลขดังต่อไปนี้

\begin{pseudocode}
\textbf{ขั้นตอนการรวมกลุ่ม:} กำหนดรายการตัวเลขจำนวนเต็มบวกตั้งแต่ 1 จำนวนขึ้นไป \\
ทำซ้ำตามขั้นตอนต่อไปนี้ \textbf{จนกว่าจะเหลือตัวเลขเพียงตัวเดียวในรายการ}
\newline
.\quad 1) ค้นหาตัวเลขที่\textbf{น้อยที่สุด 2 อันดับแรก} จากรายการ
\newline
.\quad 2) นำตัวเลขทั้งสองมารวมกัน ได้ผลรวม $S$ แล้ว\textbf{สะสม}ค่า $S$ เข้าไปใน \textbf{ต้นทุนรวม}
\newline
.\quad 3) ลบตัวเลขทั้งสองออกจากรายการ แล้วใส่ผลรวม $S$ เข้าไปแทน
\end{pseudocode}

\textbf{ตัวอย่าง:} รายการเริ่มต้น $[3,\; 1,\; 4,\; 2]$ ดำเนินการดังนี้

\begin{center}
\begin{tabular}{c|c|c|c}
\textbf{ครั้งที่} & \textbf{ตัวเลขน้อยสุด 2 อันดับ} & \textbf{ผลรวม ($S$)} & \textbf{รายการหลังรวม (เรียงใหม่)} \\
\hline
1 & $1,\; 2$ & $3$ & $[3,\; 3,\; 4]$ \\
2 & $3,\; 3$ & $6$ & $[4,\; 6]$ \\
3 & $4,\; 6$ & $10$ & $[10]$ \\
\end{tabular}
\end{center}

ต้นทุนรวม $= 3 + 6 + 10 = 19$

\textbf{คำถาม:} ถ้ารายการเริ่มต้นคือ $[2,\; 6,\; 1,\; 4,\; 3]$ เมื่อดำเนินการตามขั้นตอนข้างต้นจนเหลือตัวเลขเพียงตัวเดียว \textbf{ต้นทุนรวม} จะมีค่าเท่าใด

\choiceshorizontal{$32$}{$33$}{$34$}{$35$}

% answer: 4
% solution:
%   เริ่มต้น: [2, 6, 1, 4, 3]
%   ครั้งที่ 1: น้อยสุดคือ 1, 2 → S = 3, รายการ = [3, 6, 4, 3], ต้นทุน = 3
%   ครั้งที่ 2: น้อยสุดคือ 3, 3 → S = 6, รายการ = [4, 6, 6], ต้นทุน = 3 + 6 = 9
%   ครั้งที่ 3: น้อยสุดคือ 4, 6 → S = 10, รายการ = [6, 10], ต้นทุน = 9 + 10 = 19
%   ครั้งที่ 4: น้อยสุดคือ 6, 10 → S = 16, รายการ = [16], ต้นทุน = 19 + 16 = 35
%   ตอบ 35 (ข้อ ง)
